\subsection{Kontekst i motywacja}

\subsubsection{Wyszukiwanie informacji w bazach specjalistycznych}

Tradycyjne systemy wyszukiwania informacji opierały się na modelu boolowskim (ang. \textit{Boolean Information Retrieval}), w którym zapytania formułowane są za pomocą operatorów logicznych AND, OR oraz NOT \cite{salton1983introduction}. Podejście to, choć precyzyjne w założeniach teoretycznych, wymaga od użytkownika znajomości składni zapytań i nie uwzględnia zjawisk językowych takich jak synonimia czy polisemia \cite{manning2008introduction}. Ograniczenia te doprowadziły do rozwoju systemów typu \textit{free text retrieval}, w których zapytania formułowane są w języku naturalnym, a dokumenty rankingowane według stopnia dopasowania \cite{robertson2009probabilistic}.

Ewolucja ta jest szczególnie widoczna w kontekście specjalistycznych baz danych, takich jak PubMed czy MEDLINE, które początkowo opierały się na kontrolowanych słownikach (MeSH - \textit{Medical Subject Headings}) i zapytaniach boolowskich \cite{lu2011pubmed}. Współczesne interfejsy tych systemów oferują jednak wyszukiwanie w języku naturalnym z automatycznym rozszerzaniem zapytań o synonimy i terminy pokrewne, co znacząco obniża barierę wejścia dla użytkowników końcowych - lekarzy i badaczy nieposiadających przygotowania z zakresu informatologii. Dalszy postęp w tej dziedzinie wiąże się z zastosowaniem modeli uczenia głębokiego do reprezentacji semantycznej tekstu oraz architektur typu RAG (ang. \textit{Retrieval-Augmented Generation}), które zostały szczegółowo omówione w rozdziale 2.

\subsubsection{Rola przetwarzania mowy w wyszukiwaniu informacji}

Współczesne modele automatycznego rozpoznawania mowy (ang. \textit{Automatic Speech Recognition}, ASR) umożliwiają formułowanie zapytań do systemów wyszukiwania informacji za pomocą mowy, co jest szczególnie istotne w środowiskach klinicznych, gdzie ręczne wprowadzanie tekstu jest niepraktyczne. Modele takie jak Whisper \cite{whisper}, Wav2Vec 2.0 \cite{baevski2020wav2vec} czy HuBERT \cite{hsu2021hubert} osiągają wysoką jakość transkrypcji dla wielu języków, jednak w podejściu kaskadowym (ASR $\rightarrow$ embedding tekstowy $\rightarrow$ wyszukiwanie) błędy transkrypcji propagują się do kolejnych etapów, obniżając jakość wyników wyszukiwania. Szczegółowy przegląd architektury i możliwości tych modeli przedstawiono w rozdziale 2.

\subsubsection{Ograniczenia modeli dla języka polskiego i domen specjalistycznych}
\label{sec:motywacja}

Współczesne modele uczenia głębokiego są w przeważającej mierze trenowane na danych w języku angielskim. Zbiory danych wykorzystywane do pretrainingu dużych modeli językowych, takie jak Common Crawl, Wikipedia czy arXiv, charakteryzują się znaczącą nadreprezentacją treści anglojęzycznych \cite{joshi2020state}. Szacuje się, że język angielski stanowi ponad 90\% danych treningowych wielu popularnych modeli, mimo że jest językiem ojczystym jedynie dla około 5\% światowej populacji \cite{bender2021dangers}.

Nierównowaga ta ma istotne konsekwencje dla jakości modeli w kontekście języków innych niż angielski. Języki regionalne, w tym język polski, są klasyfikowane jako języki o średnich lub niskich zasobach (ang. \textit{mid-resource} lub \textit{low-resource languages}) \cite{joshi2020state}. Modele wielojęzyczne, takie jak mBERT \cite{devlin2019bert}, XLM-RoBERTa \cite{conneau2020unsupervised} czy Whisper \cite{whisper}, stanowią próbę rozwiązania tego problemu, jednak badania empiryczne wskazują, że ich jakość dla języków innych niż angielski jest systematycznie niższa \cite{wu2020all}, szczególnie w zadaniach wymagających znajomości specjalistycznej terminologii.

Ograniczenia modeli ogólnego przeznaczenia w kontekście domen specjalistycznych wynikają nie tylko z nierównowagi językowej, ale również z charakteru danych treningowych. Modele ASR są trenowane głównie na podcastach, filmach i nagraniach z platformy YouTube \cite{whisper}, które rzadko zawierają specjalistyczną terminologię medyczną. Analogiczny problem dotyczy modeli embeddingowych, których reprezentacje wektorowe terminów specjalistycznych mogą nie odzwierciedlać poprawnie ich znaczenia w kontekście domenowym \cite{lee2020biobert}. Istnieją rozwiązania w postaci modeli dostosowanych do konkretnych domen, jednak dotyczą one przede wszystkim języka angielskiego (zob. rozdział 2).

Kluczowym wyzwaniem pozostaje ograniczona dostępność specjalistycznych zbiorów danych dla języka polskiego. Podczas gdy dla języka angielskiego istnieją obszerne zbiory tekstów medycznych, takie jak MIMIC-III \cite{johnson2016mimic} czy PubMed, dla języka polskiego brakuje porównywalnych zasobów. Dokumentacja medyczna w języku polskim jest rozproszona, często niedostępna ze względów prawnych i wymaga pracochłonnej anonimizacji przed wykorzystaniem do celów badawczych. Sytuacja ta tworzy błędne koło: brak danych uniemożliwia trening wysokiej jakości modeli, co z kolei utrudnia automatyzację procesów, które mogłyby wspomóc tworzenie i przetwarzanie takich zbiorów.

% ----------------------------------------
\subsection{Proponowane podejście}

\subsubsection{System wyszukiwania oparty na transkrypcji audio i wektorowej bazie danych}

Klasyczna architektura systemu wyszukiwania informacji na podstawie zapytań głosowych opiera się na sekwencyjnym przetwarzaniu sygnału audio przez oddzielne moduły. Podejście to można określić jako pipeline kaskadowy, w którym każdy etap przetwarza dane wyjściowe etapu poprzedniego.

W pierwszym etapie sygnał audio jest przetwarzany przez model ASR, taki jak Whisper \cite{whisper} lub jego wersja dostrojona do domeny specjalistycznej. Wynikiem tego etapu jest transkrypcja tekstowa zapytania użytkownika. Następnie tekst jest przekształcany przez model embeddingowy w reprezentację wektorową. Modele takie jak Sentence-BERT \cite{reimers2019sentence} czy Jina Embeddings \cite{gunther2023jina} generują wektor o stałym wymiarze, reprezentujący semantyczną treść zapytania.

Uzyskany wektor zapytania jest porównywany z wektorami dokumentów zindeksowanych w wektorowej bazie danych. Systemy takie jak FAISS \cite{johnson2019billion} czy Milvus \cite{wang2021milvus} umożliwiają efektywne wyszukiwanie przybliżonych najbliższych sąsiadów (ang. \textit{approximate nearest neighbor search}) w przestrzeniach wielowymiarowych. Na podstawie miary podobieństwa system zwraca $k$ dokumentów najbardziej zbliżonych semantycznie do zapytania.

Opcjonalnie architektura może być rozszerzona o komponent generacyjny zgodnie ze wzorcem RAG \cite{lewis2020retrieval}. W takim przypadku wyszukane dokumenty stanowią kontekst dla dużego modelu językowego, który generuje odpowiedź w języku naturalnym, syntetyzując informacje z wielu źródeł. Schemat opisanej architektury przedstawiono na Rysunku \ref{fig:kaskadaArchitecture}.

\begin{figure}[htbp]
    \centering
    \begin{tikzpicture}[
    scale=0.6,
    transform shape,
    font=\small,
    block/.style={
        draw,
        rounded corners=2pt,
        thick,
        minimum width=2.8cm,
        minimum height=1.0cm,
        align=center
    },
    data/.style={block, fill=gray!12},
    model/.style={block, fill=blue!20},
    embm/.style={block, fill=purple!20},
    store/.style={block, fill=orange!25},
    ann/.style={block, fill=yellow!25},
    opt/.style={block, fill=green!15, dashed},
    arrow/.style={-{Stealth[]}, thick},
    optarrow/.style={-{Stealth[]}, thick, dashed},
    phase/.style={draw=gray, dashed, rounded corners=4pt, inner sep=0.3cm}
]

% ============================================================
% FAZA ZAPYTANIA
% ============================================================
\node[data] (audio) {Nagranie\\audio};
\node[model, right=0.7cm of audio] (asr) {Model ASR\\(Whisper)};
\node[data, right=0.7cm of asr] (txt) {Tekst\\transkrypcji};
\node[embm, right=0.7cm of txt] (qemb) {Model\\embeddingowy};
\node[data, right=1.1cm of qemb] (qvec) {Wektor\\zapytania};
\node[ann, right=1.1cm of qvec] (search) {Wyszukiwanie ANN};

\draw[arrow] (audio) -- (asr);
\draw[arrow] (asr) -- (txt);
\draw[arrow] (txt) -- (qemb);
\draw[arrow] (qemb) -- node[below, font=\scriptsize, text=gray] {$1\times d_{\mathrm{emb}}$} (qvec);
\draw[arrow] (qvec) -- (search);

% ============================================================
% FAZA INDEKSOWANIA (równolegle, offline)
% ============================================================
\node[data, above=1.6cm of qemb] (docs) {Dokumenty\\tekstowe};
\node[embm, above=1.6cm of qvec] (demb) {Model\\embeddingowy};
\node[store, above=1.6cm of search] (db) {Wektorowa\\baza danych};

\draw[arrow] (docs) -- (demb);
\draw[arrow] (demb) -- node[below, font=\scriptsize, text=gray] {$1\times d_{\mathrm{emb}}$} (db);
\draw[arrow] (db) -- node[right, font=\scriptsize] {indeks} (search);

% ============================================================
% WYNIKI + OPCJONALNY KOMPONENT GENERACYJNY (RAG)
% ============================================================
\node[data, below=0.9cm of search] (topk) {Top-$k$\\dokumentów};
\node[opt, below=0.9cm of topk] (llm) {LLM\\(generacja)};
\node[opt, below=0.9cm of llm] (ans) {Odpowiedź w\\języku naturalnym};

\draw[arrow] (search) -- (topk);
\draw[optarrow] (topk) -- node[left, font=\scriptsize\itshape] {opcjonalnie (RAG)} (llm);
\draw[optarrow] (llm) -- (ans);

% ============================================================
% RAMKI FAZ
% ============================================================
\node[phase, fit=(docs)(db),
      label={[font=\small\itshape, gray]above:Faza indeksowania (offline)}] {};
\node[phase, fit=(audio)(search),
      label={[font=\small\itshape, gray]above left:Faza zapytania}] {};

\end{tikzpicture}

    \caption{Architektura kaskadowa systemu wyszukiwania informacji na podstawie zapytań głosowych. Faza indeksowania dokumentów przebiega niezależnie od fazy zapytania. Komponent generacyjny (RAG) jest opcjonalnym rozszerzeniem architektury.}
    \label{fig:kaskadaArchitecture}
\end{figure}

Architektura kaskadowa oferuje szereg zalet praktycznych. Transkrypcja tekstowa stanowi dodatkowy produkt przetwarzania, który może być wykorzystany do celów dokumentacyjnych lub archiwalnych. Użytkownik ma możliwość weryfikacji i ewentualnej korekty rozpoznanego tekstu przed wykonaniem wyszukiwania, co zwiększa kontrolę nad procesem. Transparentność systemu ułatwia debugowanie. Każdy etap może być analizowany niezależnie, co pozwala na identyfikację źródła ewentualnych błędów. Ponadto architektura ta wykorzystuje sprawdzone, dostępne komponenty, które mogą być wymieniane niezależnie w miarę postępu technologicznego.

Podejście kaskadowe wiąże się jednak z istotnymi ograniczeniami. Dwuetapowe przetwarzanie wprowadza dodatkową latencję, co może być problematyczne w zastosowaniach wymagających odpowiedzi w czasie rzeczywistym. Błędy popełnione na etapie transkrypcji propagują się do kolejnych etapów. Niepoprawnie rozpoznany termin specjalistyczny skutkuje generacją błędnego osadzenia (ang. embedding), a w konsekwencji nieadekwatnymi wynikami wyszukiwania. 
W tym miejscu autorzy niniejszej pracy chcieliby się odnieść się do pojęcia embedding, tj. osadzenie, reprezentacja wektorowa, ale też osadzenie słów czy wektoryzacja tekstu w zależności od kontekstu, w którym się to pojęcie pojawia. Warto również zauważyć, że termin ten został spolszczony (embedding) i obecnie jest wykorzystywany w pracach związanych ze sztuczną inteligencją. Dlatgeo autorzy przyjęli, że taka właśnie forma zostanie wykorzystana w opisie prowadzonych eksperymentów.  

Ponadto redukcja sygnału audio do tekstu prowadzi do utraty informacji. Cechy prozodyczne, takie jak intonacja, tempo mowy czy wahania, mogą nieść informację o pewności lub intencji mówiącego. Informacje te są tracone w procesie transkrypcji. Przykładowo, termin medyczny wymówiony z wahaniem może wskazywać na niepewność diagnostyczną, jednak po transkrypcji jest nieodróżnialny od tego samego terminu wymówionego z pewnością.

\subsubsection{Proponowany system wyszukiwania informacji oparty na reprezentacji cech ukrytych modelu ASR}

Alternatywnym podejściem jest bezpośrednie wykorzystanie reprezentacji ukrytych generowanych przez enkoder modelu ASR do wyszukiwania informacji, z pominięciem etapu transkrypcji tekstowej. Koncepcja ta, określana jako Speech-to-Retrieval, opiera się na założeniu, że reprezentacje ukryte modeli ASR zawierają informację semantyczną wystarczającą do realizacji zadań wyszukiwania.

Modele takie jak Wav2Vec 2.0 \cite{baevski2020wav2vec}, HuBERT \cite{hsu2021hubert} czy enkoder modelu Whisper generują w warstwach pośrednich reprezentacje wektorowe segmentów audio. Reprezentacje te, będące wynikiem przetwarzania przez głębokie sieci neuronowe, kodują informację na różnych poziomach abstrakcji. Warstwy bliższe wejściu reprezentują cechy akustyczne niskiego poziomu, podczas gdy warstwy głębsze mogą kodować informację fonologiczną, leksykalną, a potencjalnie również semantyczną \cite{pasad2021layer}.

Architektura proponowanego systemu składa się z enkodera ASR, warstwy projekcji oraz wektorowej bazy danych. Sygnał audio jest przetwarzany przez enkoder, który generuje sekwencję reprezentacji ukrytych. Sekwencja ta jest następnie agregowana (np. przez uśrednienie lub mechanizm uwagi) i przekształcana przez warstwę projekcji do przestrzeni wektorowej zgodnej z embeddingami tekstowymi dokumentów. Uzyskany wektor służy do wyszukiwania w bazie wektorowej analogicznie jak w architekturze kaskadowej. Schemat proponowanej architektury przedstawiono na Rysunku \ref{fig:s2rArchitecture}.

\begin{figure}[htbp]
    \centering
    \begin{tikzpicture}[
    scale=0.55,
    transform shape,
    font=\small,
    block/.style={
        draw,
        rounded corners=2pt,
        thick,
        minimum width=2.6cm,
        minimum height=1.0cm,
        align=center
    },
    data/.style={block, fill=gray!12},
    model/.style={block, fill=blue!20},
    embm/.style={block, fill=purple!20},
    adapt/.style={block, fill=teal!25},
    store/.style={block, fill=orange!25},
    ann/.style={block, fill=yellow!25},
    opt/.style={block, fill=green!15, dashed},
    arrow/.style={-{Stealth[]}, thick},
    optarrow/.style={-{Stealth[]}, thick, dashed},
    phase/.style={draw=gray, dashed, rounded corners=4pt, inner sep=0.4cm}
]

% ============================================================
% FAZA ZAPYTANIA (bez etapu transkrypcji)
% ============================================================
\node[data] (audio) {Nagranie\\audio};
\node[model, right=0.9cm of audio] (enc) {Enkoder ASR};
\node[data, right=1.5cm of enc] (latent) {Reprezentacja\\ukryta};
\node[model, right=0.9cm of latent] (agg) {Agregacja\\sekwencji};
\node[adapt, right=1.5cm of agg] (proj) {Warstwa projekcji\\(adapter)};
\node[data, right=1.9cm of proj] (qvec) {Wektor\\zapytania};
\node[ann, right=1.9cm of qvec] (search) {Wyszukiwanie ANN};

\draw[arrow] (audio) -- (enc);
\draw[arrow] (enc) -- node[above, font=\scriptsize] {$h_{1..T}$}
    node[below, font=\scriptsize, text=gray] {$T\times d$} (latent);
\draw[arrow] (latent) -- (agg);
\draw[arrow] (agg) -- node[above, font=\scriptsize] {$\bar{h}$}
    node[below, font=\scriptsize, text=gray] {$1\times d$} (proj);
\draw[arrow] (proj) -- node[below, font=\scriptsize, text=gray] {$1\times d_{\mathrm{emb}}$} (qvec);
\draw[arrow] (qvec) -- (search);

% ============================================================
% FAZA INDEKSOWANIA (dokumenty nadal indeksowane tekstowo)
% ============================================================
\node[data, above=2.6cm of proj] (docs) {Dokumenty\\tekstowe};
\node[embm, above=2.6cm of qvec] (demb) {Model\\embeddingowy};
\node[store, above=2.6cm of search] (db) {Wektorowa\\baza danych};

\draw[arrow] (docs) -- (demb);
\draw[arrow] (demb) -- node[below, font=\scriptsize, text=gray] {$1\times d_{\mathrm{emb}}$} (db);
\draw[arrow] (db) -- node[right, font=\scriptsize] {indeks} (search);

% wspólna przestrzeń wektorowa: cel treningu adaptera
\draw[{Stealth[]}-{Stealth[]}, dotted, thick, gray]
    (demb.south) -- node[left, font=\scriptsize\itshape, text=gray]
    {wspólna przestrzeń} (qvec.north);

% ============================================================
% WYNIKI + OPCJONALNY KOMPONENT GENERACYJNY (RAG)
% ============================================================
\node[data, below=1.2cm of search] (topk) {Top-$k$\\dokumentów};
\node[opt, below=1.2cm of topk] (llm) {LLM\\(generacja)};
\node[opt, below=1.2cm of llm] (ans) {Odpowiedź w\\języku naturalnym};

\draw[arrow] (search) -- (topk);
\draw[optarrow] (topk) -- node[left, font=\scriptsize\itshape] {opcjonalnie (RAG)} (llm);
\draw[optarrow] (llm) -- (ans);

% ============================================================
% RAMKI FAZ
% ============================================================
\node[phase, fit=(docs)(db),
      label={[font=\small\itshape, gray]above:Faza indeksowania (offline)}] {};
\node[phase, fit=(audio)(search),
      label={[font=\small\itshape, gray]above left:Faza zapytania (bez transkrypcji)}] {};

\end{tikzpicture}

    \caption{Architektura systemu wyszukiwania opartego na reprezentacji cech ukrytych modelu ASR (\textit{Speech-to-Retrieval}). Etap transkrypcji tekstowej zostaje pominięty, a warstwa projekcji odwzorowuje reprezentację audio do przestrzeni wektorowej modelu embeddingowego, w której zindeksowane są dokumenty tekstowe.}
    \label{fig:s2rArchitecture}
\end{figure}

Kluczowymi elementami architektury są agregacja sekwencji oraz warstwa projekcji (adapter). Zadaniem agregacji jest wydobycie najistotniejszych informacji z sekwencji reprezentacji ukrytych i sprowadzenie jej do pojedynczego wektora, czyli spłaszczenie reprezentacji o wymiar czasu ($T \times d \rightarrow 1 \times d$). Warstwa projekcji odwzorowuje następnie tak uzyskaną reprezentację z przestrzeni enkodera ASR do przestrzeni embeddingów tekstowych. Trening tych komponentów wymaga zbioru par (audio, tekst) wraz z odpowiadającymi im embeddingami tekstowymi. Celem optymalizacji jest minimalizacja odległości między projekcją reprezentacji audio a embeddingiem odpowiadającego tekstu.

Potencjalne zalety proponowanego podejścia obejmują redukcję latencji poprzez eliminację etapu dekodowania tekstu oraz możliwość wykorzystania informacji akustycznej niedostępnej w transkrypcji. Ponadto architektura end-to-end otwiera możliwość wspólnej optymalizacji komponentów retrieval i ASR, co może prowadzić do lepszego dostosowania reprezentacji do zadania wyszukiwania.

Podejście to wiąże się jednak z istotnymi wyzwaniami. Brak transkrypcji utrudnia weryfikację i interpretację działania systemu przez użytkownika. Wyjaśnialność systemu jest ograniczona, ponieważ trudniej jest zrozumieć, dlaczego system zwrócił określone dokumenty. Trening warstwy projekcji wymaga odpowiednio dużego zbioru danych. Ponadto modele ASR nie były pierwotnie trenowane z myślą o zadaniach wyszukiwania informacji, co może ograniczać jakość reprezentacji semantycznych.

\subsubsection{Pytania badawcze}

Porównanie przedstawionych architektur prowadzi do sformułowania pytań badawczych stanowiących przedmiot niniejszej pracy:

\begin{enumerate}
    \item Czy reprezentacje ukryte generowane przez modele ASR zawierają wystarczającą informację semantyczną do realizacji zadań wyszukiwania informacji?
    \item Jak jakość wyszukiwania opartego na reprezentacjach ukrytych wypada w porównaniu z podejściem kaskadowym (ASR → embedding tekstowy)?
    \item Jaki wpływ na jakość obu podejść ma język zapytania (polski vs. angielski) oraz domena (terminologia medyczna vs. język ogólny)?
    \item Które warstwy modeli ASR generują reprezentacje najbardziej przydatne do zadań wyszukiwania semantycznego?
\end{enumerate}

Odpowiedzi na te pytania pozwolą na ocenę praktycznej przydatności podejścia Speech-to-Retrieval w kontekście specjalistycznych systemów wyszukiwania informacji dla języka polskiego.

% ----------------------------------------
\subsection{Struktura pracy}

Dalsza część pracy jest zorganizowana w następujący sposób. Rozdział 2 zawiera przegląd literatury obejmujący systemy wyszukiwania informacji, modele embeddingowe,
architekturę RAG oraz współczesne modele ASR. W rozdziale 3 przedstawiono badania nad semantycznością reprezentacji wektorów cech modeli ASR.
Rozdział 4 opisuje ewaluację architektur dopasowania reprezentacji latentnych modeli ASR i modeli embeddingowych. 
Rozdział 5 dzieli się na dwie części. Pierwsza przedstawia system ewaluacji zbudowany na potrzeby pracy. Druga zawiera eksperymenty porównujące architekturę kaskadową
z proponowanym dopasowaniem reprezentacji latentnych oraz omówienie uzyskanych wyników.
Rozdział 6 zawiera podsumowanie i wnioski.
